
\section{\BookRecordIndex{T0184}简单稳定映射的\BookRecordIndex{T0166}横截性}\label{sec:6-2}
\zhindex{jiandanwendingyingshedehengjiexing}{简单稳定映射的横截性}

固定一棵 $k$-\BookRecordIndex{T0084}带标号树 $T=(T,E,\Lambda)$，以及满足稳定性条件
\eqref{eq:6-1-1} 的函数
\[
  T\longrightarrow H_2(M;\mathbb Z),
  \qquad \alpha\longmapsto A_\alpha.
\]
记
\[
  A:=\sum_{\alpha\in T}A_\alpha.
\]
我们要证明：对通有的\BookRecordIndex{T0088}殆复结构 $J\in\mathcal J_\tau(M,\omega)$，\BookRecordIndex{T0246}模空间
$\mathcal M^*_{0,T}(\{A_\alpha\};J)$ 是一个具有“正确”维数的光滑、定向有限维流形。
第一步是表述为此所需的 $J$ 应满足的条件。显然，对每个类 $A_\alpha$，$J$ 都应当在
定义~\ref{def:3-1-5} 的意义下正则，从而每个\BookRecordIndex{T0246}模空间
$\mathcal M^*(A_\alpha;J)$ 都是流形。然而，这还不够：还需要不同空间
$\mathcal M^*(A_\alpha;J)$ 彼此横截。

为理解这里需要什么条件，设 $\{A_\alpha\}_{\alpha\in T}$ 是
$H_2(M;\mathbb Z)$ 中任意有限个球面同调类，考虑由 $J$-全纯球面元组
$\mathbf u=\{u_\alpha\}_{\alpha\in T}$ 所成的空间
\[
  \mathcal M^*(\{A_\alpha\};J)
  \subset\prod_{\alpha\in T}\mathcal M^*(A_\alpha;J),
\]
其中 $[u_\alpha]=A_\alpha$，每个非常值 $u_\alpha$ 都单纯，并且只要
$\alpha\ne\beta$ 且 $u_\alpha,u_\beta$ 都非常值，就有
$u_\alpha(S^2)\ne u_\beta(S^2)$。因此，
$\mathcal M^*(\{A_\alpha\};J)$ 是第~\ref{sec:3-4} 节记作
$\mathcal M^*(A,\Sigma;J)$ 的、定义域
$\Sigma:=T\times S^2$ 不连通的\BookRecordIndex{T0091}单纯曲线空间的一个子集。

再考虑开集
\[
  Z(T)\subset(S^2)^E\times(S^2)^k,
\]
它由所有元组
\[
  \mathbf z=\bigl(\{z_{\alpha\beta}\}_{\alpha E\beta},
  \{z_i\}_{1\le i\le k}\bigr)
\]
组成，并且对每个 $\alpha$，所有满足 $\alpha E\beta$ 的点 $z_{\alpha\beta}$ 与所有满足
$\alpha_i=\alpha$ 的点 $z_i$ 两两不同。这里可以把 $E\subset T\times T$ 的元素看作
有向边。不过，如附录 D.2 所述，我们的一般约定是：除非明确另作说明，边都是无向的。
有一个\emph{\BookRecordIndex{T0065}边求值映射}
\[
  \operatorname{ev}^E:
  \mathcal M^*(\{A_\alpha\};J)\times Z(T)
  \longrightarrow M^E,
\]
定义为
\begin{equation}\label{eq:6-2-1}
  \operatorname{ev}^E(\mathbf u,\mathbf z)
  :=\{u_\alpha(z_{\alpha\beta})\}_{\alpha E\beta}.
\end{equation}
记 $\Delta^E\subset M^E$ 为由边关系确定的对角线：
\[
  \Delta^E
  :=\bigl\{\{x_{\alpha\beta}\}_{\alpha E\beta}\in M^E
  \bigm|x_{\alpha\beta}=x_{\beta\alpha}\bigr\}.
\]

\begin{definition}\label{def:6-2-1}
若被 $\omega$ \BookRecordIndex{T0372}驯服的\BookRecordIndex{T0088}殆复结构 $J\in\mathcal J_\tau(M,\omega)$ 满足下列条件，就称它对
$T$ 和 $\{A_\alpha\}$ \emph{正则}。
\begin{enumerate}[label=(\roman*),leftmargin=3em]
\item 对每个 $\mathbf u=\{u_\alpha\}\in
  \mathcal M^*(\{A_\alpha\};J)$ 和每个 $\alpha\in T$，$D_{u_\alpha}$ 都满射。
\item \BookRecordIndex{T0283}求值映射 $\operatorname{ev}^E$ 与 $\Delta^E$ 横截。
\end{enumerate}
\end{definition}

注意，条件 (i) 蕴含 $\operatorname{ev}^E$ 的定义域
$\mathcal M^*(\{A_\alpha\};J)\times Z(T)$ 是流形，所以条件 (ii) 是有意义的。下面给出
一个更显式的版本。

\begin{lemma}\label{lem:6-2-2}
\BookRecordIndex{T0065}边求值映射 $\operatorname{ev}^E$ 与 $\Delta^E$ 横截，当且仅当对每个\BookRecordIndex{T0184}简单稳定映射
\[
  (\mathbf u,\mathbf z)
  =\bigl(\{u_\alpha\}_{\alpha\in T},
    \{z_{\alpha\beta}\}_{\alpha E\beta},
    \{z_i\}_{1\le i\le k}\bigr)
  \in\widetilde{\mathcal M}^*_{0,T}(\{A_\alpha\};J)
\]
和每族满足
\[
  v_{\alpha\beta}\in T_{u_\alpha(z_{\alpha\beta})}M,
  \qquad v_{\alpha\beta}+v_{\beta\alpha}=0
\]
的切向量，都存在向量场 $\xi_\alpha\in\ker D_{u_\alpha}$ 和切向量
$\zeta_{\alpha\beta}\in T_{z_{\alpha\beta}}S^2$，使得对所有满足
$\alpha E\beta$ 的 $\alpha,\beta\in T$，
\begin{equation}\label{eq:6-2-2}
  v_{\alpha\beta}
  =\xi_\alpha(z_{\alpha\beta})
   +du_\alpha(z_{\alpha\beta})\zeta_{\alpha\beta}
   -\xi_\beta(z_{\beta\alpha})
   -du_\beta(z_{\beta\alpha})\zeta_{\beta\alpha}.
\end{equation}
\end{lemma}

\begin{proof}
把\BookRecordIndex{T0065}边求值映射写成
\[
  \operatorname{ev}^E
  =\{\operatorname{ev}_{\alpha\beta}\}_{\alpha E\beta},
\]
其中
\[
  \operatorname{ev}_{\alpha\beta}:
  \mathcal M^*(\{A_\alpha\};J)\times Z(T)\longrightarrow M,
  \qquad
  \operatorname{ev}_{\alpha\beta}(\mathbf u,\mathbf z)
  :=u_\alpha(z_{\alpha\beta}).
\]
于是，对 $\boldsymbol\xi=\{\xi_\alpha\}_{\alpha\in T}$ 和
$\boldsymbol\zeta=\{\zeta_{\alpha\beta}\}_{\alpha E\beta}$，
\[
  d\operatorname{ev}_{\alpha\beta}(\mathbf u,\mathbf z)
  (\boldsymbol\xi,\boldsymbol\zeta)
  =\xi_\alpha(z_{\alpha\beta})
   +du_\alpha(z_{\alpha\beta})\zeta_{\alpha\beta}.
\]
因此，式~\eqref{eq:6-2-2} 等价于
\[
  v_{\alpha\beta}
  =d\operatorname{ev}_{\alpha\beta}(\mathbf u,\mathbf z)
    (\boldsymbol\xi,\boldsymbol\zeta)
   -d\operatorname{ev}_{\beta\alpha}(\mathbf u,\mathbf z)
    (\boldsymbol\xi,\boldsymbol\zeta)
\]
对 $\alpha E\beta$ 成立。作此理解后，结论由习题~\ref{ex:6-2-3} 得出。
\end{proof}

\begin{exercise}\label{ex:6-2-3}
设 $V$ 是向量空间，$(T,E)$ 是一棵树，$\Delta^E\subset V^E$ 是\BookRecordIndex{T0061}边对角线，
$W\subset V^E$ 是一个子空间。证明
\[
  V^E=\Delta^E+W
\]
当且仅当对每个满足 $v_{\alpha\beta}+v_{\beta\alpha}=0$ 的元组
$\{v_{\alpha\beta}\}_{\alpha E\beta}\in V^E$，都存在元组
$\{w_{\alpha\beta}\}_{\alpha E\beta}\in W$，使得对 $\alpha E\beta$，
$v_{\alpha\beta}=w_{\alpha\beta}-w_{\beta\alpha}$。
\end{exercise}

\begin{exercise}\label{ex:6-2-4}
设 $\mathcal M^*(\{A_\alpha\};J)\ne\varnothing$。证明定义~\ref{def:6-2-1} 的条件
(i) 成立，当且仅当 $J$ 对所有 $A_\alpha$ 都在定义~\ref{def:3-1-5} 的意义下正则。
注意，只要存在一对顶点 $\alpha,\beta\in T$，满足
$\alpha\ne\beta$、$A_\alpha=A_\beta$，且类 $A_\alpha$ 中只有一个 $J$-全纯球面
（不计重参数化），就有 $\mathcal M^*(\{A_\alpha\};J)=\varnothing$。在什么其他情形下，
简单元组的\BookRecordIndex{T0246}模空间为空？
\end{exercise}

下一个习题说明，添加\BookRecordIndex{T0066}标记点不影响\BookRecordIndex{T0088}殆复结构的\BookRecordIndex{T0408}正则性。

\begin{exercise}\label{ex:6-2-5}
设 $T=(T,E,\Lambda)$ 是一棵 $k$-\BookRecordIndex{T0084}带标号树，并且映射
$T\to H_2(M;\mathbb Z):\alpha\mapsto A_\alpha$ 满足稳定性条件
\eqref{eq:6-1-1}。遗忘\BookRecordIndex{T0066}标记点，并删去所有满足 $A_\alpha=0$ 且承载少于三个\BookRecordIndex{T0191}节点的顶点，
由此得到约化树 $(T',E')$。证明：$J$ 对 $T$ 和
$\{A_\alpha\}_{\alpha\in T}$ 正则，当且仅当它对 $T'$ 和
$\{A_\alpha\}_{\alpha\in T'}$ 正则。提示：证明只需考虑
$T=T'\cup\{\gamma\}$、且至多有两条边以 $\gamma$ 为端点的情形。利用如下事实：当
$u_\alpha$ 是常值映射时，$\ker D_{u_\alpha}$ 由平凡丛
$u_\alpha^*TM=S^2\times T_{u_\alpha(z)}M$ 的常值截面组成。定理~\ref{thm:6-3-1}
在第~\ref{sec:6-3} 节给出的证明中也会用到相关论证。
\end{exercise}

以下用 $\mathcal J_{\mathrm{reg}}(T,\{A_\alpha\})$ 表示所有对
$T$ 和 $\{A_\alpha\}$ 正则的被 $\omega$ \BookRecordIndex{T0372}驯服\BookRecordIndex{T0088}殆复结构所成的集合，并记
\[
  \mathcal J_{\mathrm{reg}}(M,\omega)
  :=\bigcap_T\bigcap_{\{A_\alpha\}}
  \mathcal J_{\mathrm{reg}}(T,\{A_\alpha\})
\]
为对每棵 $T$ 和每个 $\{A_\alpha\}$ 都正则的\BookRecordIndex{T0088}殆复结构集合。注意这是可数交。一个\BookRecordIndex{T0088}殆复结构
$J\in\mathcal J_\tau(M,\omega)$ 属于 $\mathcal J_{\mathrm{reg}}(M,\omega)$，当且仅当
对每个单纯 $J$-全纯球面 $u$，算子 $D_u$ 都满射，并且对每棵 $T$ 与每个
$\{A_\alpha\}$，\BookRecordIndex{T0283}求值映射~\eqref{eq:6-2-1} 都与 $\Delta^E$ 横截。准备至此，可以陈述本节
的主要定理了。

\begin{theorem}\label{thm:6-2-6}
设 $T$ 是一棵 $k$-\BookRecordIndex{T0084}带标号树，并固定满足稳定性条件~\eqref{eq:6-1-1} 的映射
\[
  T\longrightarrow H_2(M;\mathbb Z),
  \qquad \alpha\longmapsto A_\alpha.
\]
\begin{enumerate}[label=(\roman*),leftmargin=3em]
\item 若 $J\in\mathcal J_{\mathrm{reg}}(T,\{A_\alpha\})$，则\BookRecordIndex{T0246}模空间
$\mathcal M^*_{0,T}(\{A_\alpha\};J)$ 是光滑定向流形，其维数为
\[
  \dim\mathcal M^*_{0,T}(\{A_\alpha\};J)
  =\mu(A,T):=2n+2c_1(A)+2k(T)-6-2e(T),
\]
其中 $k(T)$ 是\BookRecordIndex{T0066}标记点数，$e(T)$ 是边数。
\item 集合 $\mathcal J_{\mathrm{reg}}(T,\{A_\alpha\})$ 是
$\mathcal J_\tau(M,\omega)$ 中的\BookRecordIndex{T0300}剩余集。
\end{enumerate}
\end{theorem}

\begin{proof}[定理~\ref{thm:6-2-6} 的 (i)]
固定 $J\in\mathcal J_{\mathrm{reg}}(T,\{A_\alpha\})$。由定义~\ref{def:6-2-1}
的 (i) 与定理~\ref{thm:3-1-6}，\BookRecordIndex{T0246}模空间
$\mathcal M^*(\{A_\alpha\};J)$ 是光滑定向流形，维数为
\[
  \dim\mathcal M^*(\{A_\alpha\};J)
  =2n\bigl(e(T)+1\bigr)+2c_1(A).
\]
由定义~\ref{def:6-2-1} 的 (ii)，\BookRecordIndex{T0283}求值映射
\[
  \operatorname{ev}^E:
  \mathcal M^*(\{A_\alpha\};J)\times Z(T)\longrightarrow M^E
\]
与 $\Delta^E$ 横截，因此原像
\[
  \widetilde{\mathcal M}^*_{0,T}(\{A_\alpha\};J)
  =\bigl(\operatorname{ev}^E\bigr)^{-1}(\Delta^E)
\]
\sourcefootnote{原书在此处及第~\ref{sec:6-3} 节开头两次把原像映射误印为
$\pi^E$；本章定义的相应映射是\BookRecordIndex{T0065}边求值映射 $\operatorname{ev}^E$，且原像条件正是
\BookRecordIndex{T0191}节点两端取值相等，故两处均改正。}
是光滑定向流形。由于 $Z(T)$ 的维数为 $4e(T)+2k(T)$，而 $\Delta^E$ 的余维数为
$2ne(T)$，故
\[
  \dim\widetilde{\mathcal M}^*_{0,T}(\{A_\alpha\};J)
  =2n+2c_1(A)+2k(T)+4e(T).
\]
又因为 $G_T$ 的维数为 $6e(T)+6$，并且它在
$\widetilde{\mathcal M}^*_{0,T}(\{A_\alpha\};J)$
\sourcefootnote{原书此处把作用对象印作
$\widetilde{\mathcal M}^{*}(\{A_\alpha\};J)$，漏掉下标 $0,T$；前一显示刚定义的原像空间
才含有\BookRecordIndex{T0191}节点数据、其商才是本定理中的\BookRecordIndex{T0340}稳定映射\BookRecordIndex{T0246}模空间，故补回下标。}
上通过保定向微分同胚自由且恰当地作用，这就证明了定理~\ref{thm:6-2-6} 的 (i)。
\end{proof}

(ii) 的证明建立在命题~\ref{prop:3-2-1} 的如下推广之上。设 $T$ 是任意有限集，
$\{A_\alpha\}_{\alpha\in T}$ 是任意一族球面同调类。记
$\mathcal J^\ell=\mathcal J^\ell_\tau(M,\omega)$ 为所有 $C^\ell$ 类、被 $\omega$ \BookRecordIndex{T0372}驯服的
\BookRecordIndex{T0088}殆复结构所成的\BookRecordIndex{T0053}巴拿赫流形。相应的\BookRecordIndex{T0268}普适模空间为
\[
  \mathcal M^*(\{A_\alpha\};\mathcal J^\ell)
  :=\bigl\{(\mathbf u,J)\bigm|
    J\in\mathcal J^\ell,
    \ \mathbf u\in\mathcal M^*(\{A_\alpha\};J)\bigr\}.
\]
与第~\ref{sec:3-4} 节一样，可以把这个空间看作\BookRecordIndex{T0053}巴拿赫流形
$\mathcal B^{k,p}\times\mathcal J^\ell$ 的子集，其中
\[
  \mathcal B^{k,p}
  :=\bigl\{\{u_\alpha\}_{\alpha\in T}\bigm|
  u_\alpha\in W^{k,p}(S^2,M),\ [u_\alpha]=A_\alpha\bigr\}.
\]

\begin{proposition}\label{prop:6-2-7}
设 $\ell\ge2$ 是整数，$p>2$ 是实数，
$\{A_\alpha\}_{\alpha\in T}$ 是 $H_2(M;\mathbb Z)$ 中有限个球面同调类，并且
$k\in\{1,\ldots,\ell\}$。则
$\mathcal M^*(\{A_\alpha\}_{\alpha\in T};\mathcal J^\ell)$ 是
$\mathcal B^{k,p}\times\mathcal J^\ell$ 的一个 $C^{\ell-k}$ 巴拿赫子流形。
\end{proposition}

\begin{proof}
证明几乎逐字照搬命题~\ref{prop:3-2-1} 的证明。唯一的区别是，在\BookRecordIndex{T0166}横截性证明中，我们必须
同时处理一族 $L^q$ 类 $(0,1)$-形式
$\eta_\alpha\in\Omega^{0,1}(S^2,u_\alpha^*TM)$，并须证明：若对每个 $\alpha\in T$ 都有
$D_{u_\alpha}^*\eta_\alpha=0$，而且对每个
$Y\in T_J\mathcal J^\ell$ 都有
\begin{equation}\label{eq:6-2-3}
  \sum_{\alpha\in T}\int_{S^2}
  \bigl\langle\eta_\alpha,
  Y(u_\alpha)du_\alpha\circ i\bigr\rangle
  \,d\operatorname{vol}_{S^2}=0,
\end{equation}
则所有 $\eta_\alpha$ 都为零。

接下来的证明与命题~\ref{prop:3-2-1} 相同，只是在选取 $Y$ 时必须更加谨慎。令
$Z(u_\alpha)\subset S^2$ 表示 $u_\alpha$ 的\BookRecordIndex{T0119}非单射点集。由命题~\ref{prop:2-5-1}，
$Z(u_\alpha)$ 至多可数，并且只能在 $u_\alpha$ 的\BookRecordIndex{T0233}临界点处聚集。由
命题~\ref{prop:2-4-4}，当 $\beta\ne\alpha$ 时，集合
$u_\alpha^{-1}(u_\beta(S^2))$ 也有相同性质。于是，仿照命题~\ref{prop:3-2-1} 的证明，
由式~\eqref{eq:6-2-3} 可知，$\eta_\alpha$ 必须在
\[
  Z(u_\alpha)\cup
  \bigcup_{\beta\ne\alpha}u_\alpha^{-1}(u_\beta(S^2))
\]
的补集中的每一点都为零。该补集在 $S^2$ 中开且稠密，故对每个 $\alpha$ 都有
$\eta_\alpha\equiv0$。作此理解后，论证便与命题~\ref{prop:3-2-1} 的证明相同。这证明了
命题~\ref{prop:6-2-7}。
\end{proof}

现在固定一棵 $k$-\BookRecordIndex{T0084}带标号树 $T=(T,E,\Lambda)$，考虑普适\BookRecordIndex{T0065}边求值映射
\[
  \operatorname{ev}^E:
  \mathcal M^*(\{A_\alpha\};\mathcal J^\ell)\times Z(T)
  \longrightarrow M^E,
\]
定义为
\[
  \operatorname{ev}^E(\mathbf u,\mathbf z,J)
  :=\{u_\alpha(z_{\alpha\beta})\}_{\alpha E\beta}.
\]

\begin{proposition}\label{prop:6-2-8}
对每个 $\ell\ge2$、每棵 $k$-\BookRecordIndex{T0084}带标号树 $T=(T,E,\Lambda)$，以及
$H_2(M;\mathbb Z)$ 中每族球面同调类 $\{A_\alpha\}_{\alpha\in T}$，普适\BookRecordIndex{T0065}边求值映射
$\operatorname{ev}^E$ 都与 $\Delta^E$ 横截。
\end{proposition}

\begin{proof}
见第~\ref{sec:6-3} 节。
\end{proof}

当 $T$ 只有一条边时，$\Delta^E$ 是 $M^2$ 的子流形，因此该命题由
命题~\ref{prop:3-4-2} 得出。一般而言，对固定的 $E$，$\Delta^E$ 是各条边所对应的
对角子流形之交；等价地，它是各成对子对角线的乘积。举例说，若有两条边，那么
$\Delta^E$ 是 $M^4$ 中分别由方程 $x_1=x_2$ 与 $x_3=x_4$ 给出的两个子流形之交
\sourcefootnote{原书把 $\Delta^E$ 两度描述为这些子流形的“并”，并进而称其为由不同维数
子流形组成的分\BookRecordIndex{T0072}层空间；但定义~\ref{def:6-2-1} 要求同时满足所有边等式，且下文使用的
$\operatorname{codim}\Delta^E=2ne(T)$ 也只与交（等价于成对子对角线的乘积）一致，故据此
改正。}。推广到多条边并不困难——我们将对边数作归纳——不过，为免打断当前论证的行文，
把它推迟到下一节。

\begin{proof}[定理~\ref{thm:6-2-6} 的 (ii)]
由命题~\ref{prop:6-2-7}，
$\mathcal M^*(\{A_\alpha\};\mathcal J^\ell)$ 是\BookRecordIndex{T0213}可分的
$C^{\ell-1}$ \BookRecordIndex{T0053}巴拿赫流形；由命题~\ref{prop:6-2-8}，\BookRecordIndex{T0269}普适求值映射
$\operatorname{ev}^E$ 与 $\Delta^E$ 横截。因此，由\BookRecordIndex{T0380}隐函数定理 A.3.3，
$\Delta^E$ 在 $\operatorname{ev}^E$ 下的原像是\BookRecordIndex{T0213}可分的 $C^{\ell-1}$ \BookRecordIndex{T0053}巴拿赫流形。这个原像
就是以树 $T$ 为模型、表示同调类 $A_\alpha$、对应某个
$J\in\mathcal J^\ell$ 的\BookRecordIndex{T0184}简单稳定映射的\BookRecordIndex{T0268}普适模空间
\begin{equation}\label{eq:6-2-4}
  \widetilde{\mathcal M}^*_{0,T}(\{A_\alpha\};\mathcal J^\ell)
  :=\bigl\{(\mathbf u,\mathbf z,J)\bigm|
  J\in\mathcal J^\ell,
  \ (\mathbf u,\mathbf z)\in
  \widetilde{\mathcal M}^*_{0,T}(\{A_\alpha\};J)\bigr\}.
\end{equation}
现在考虑投影
\[
  p^\ell:\mathcal M^*(\{A_\alpha\};\mathcal J^\ell)
  \longrightarrow\mathcal J^\ell,
  \qquad
  \pi^\ell:\widetilde{\mathcal M}^*_{0,T}
  (\{A_\alpha\};\mathcal J^\ell)\longrightarrow\mathcal J^\ell.
\]
它们都是弗雷德霍姆映射，指标分别为
\[
  \operatorname{index}(p^\ell)
  =2n\bigl(1+e(T)\bigr)+2c_1(A),
  \qquad
  \operatorname{index}(\pi^\ell)=\mu(A,T)+\dim G_T.
\]
于是，由 萨德--斯梅尔定理 A.5.1，当 $\ell$ 充分大时，$p^\ell$ 与 $\pi^\ell$ 的共同
\BookRecordIndex{T0409}正则值集合
$\mathcal J^\ell_{\mathrm{reg}}(T,\{A_\alpha\})$ 是
$\mathcal J^\ell$ 中的\BookRecordIndex{T0300}剩余集。此外，\BookRecordIndex{T0088}殆复结构 $J\in\mathcal J^\ell$ 是
$p^\ell$ 和 $\pi^\ell$ 的共同\BookRecordIndex{T0409}正则值，当且仅当它满足定义~\ref{def:6-2-1} 的条件。

现在，对每个 $K>0$，考虑子集
$\mathcal M^*_K(\{A_\alpha\};\mathcal J^\ell)
\subset\mathcal M^*(\{A_\alpha\};\mathcal J^\ell)$，它由所有满足下列条件的元组
$(\mathbf u,J)$ 组成：对每个 $\alpha\in T$，
\begin{equation}\label{eq:6-2-5}
  \lVert du_\alpha\rVert_{L^\infty}\le K,
\end{equation}
而且存在一族点 $\{z_\alpha\}_{\alpha\in T}\subset S^2$，使得对每个
$\alpha\in T$ 及 $\beta\in T\setminus\{\alpha\}$，
\begin{equation}\label{eq:6-2-6}
  \inf_{\zeta\in S^2\setminus\{z_\alpha\}}
  \frac{d\bigl(u_\alpha(z_\alpha),u_\alpha(\zeta)\bigr)}
       {d(z_\alpha,\zeta)}
  \ge\frac1K,
  \qquad
  \inf_{\zeta\in S^2}
  d\bigl(u_\alpha(z_\alpha),u_\beta(\zeta)\bigr)
  \ge\frac1K.
\end{equation}
类似地，令 $Z_K(T)\subset Z(T)$ 为所有满足如下条件的元组
$\mathbf z\in Z(T)$ 所成的集合：只要相应各点位于同一个 $\alpha$-球面上，就有
\[
  d(z_{\alpha\beta},z_{\alpha\gamma})\ge\frac1K,
  \qquad
  d(z_i,z_j)\ge\frac1K,
  \qquad
  d(z_{\alpha\beta},z_i)\ge\frac1K;
\]
更明确地说，这些不等式分别对
$\alpha E\beta$、$\alpha E\gamma$、$\beta\ne\gamma$，
$i\ne j$、$\alpha_i=\alpha_j=\alpha$，以及
$\alpha E\beta$、$\alpha_i=\alpha$ 成立。对 $J\in\mathcal J^\ell$，记
\[
  \widetilde{\mathcal M}^*_{0,T;K}(\{A_\alpha\};J)
  :=\bigl\{(\mathbf u,\mathbf z)\in
    \widetilde{\mathcal M}^*_{0,T}(\{A_\alpha\};J)
    \bigm|(\mathbf u,J)\in
    \mathcal M^*_K(\{A_\alpha\};\mathcal J^\ell),\quad
    \mathbf z\in Z_K(T)\bigr\},
\]
并以
$\widetilde{\mathcal M}^*_{0,T;K}
(\{A_\alpha\};\mathcal J^\ell)$
表示这些纤维组成的普适空间。
于是，投影
\[
  p_K^\ell:\mathcal M^*_K(\{A_\alpha\};\mathcal J^\ell)
  \longrightarrow\mathcal J^\ell,
  \qquad
  \pi_K^\ell:\widetilde{\mathcal M}^*_{0,T;K}
  (\{A_\alpha\};\mathcal J^\ell)\longrightarrow\mathcal J^\ell
\]
是恰当弗雷德霍姆映射。因此，当 $\ell$ 充分大时，$p_K^\ell$ 与 $\pi_K^\ell$ 的共同
\BookRecordIndex{T0409}正则值集合 $\mathcal J^\ell_{\mathrm{reg},K}(T,\{A_\alpha\})$ 在
$\mathcal J^\ell$ 中开且稠密。同理，集合
\[
  \mathcal J_{\mathrm{reg},K}(T,\{A_\alpha\})
  :=\mathcal J^\ell_{\mathrm{reg},K}(T,\{A_\alpha\})
  \cap\mathcal J_\tau(M,\omega)
\]
在 $\mathcal J_\tau(M,\omega)$ 中开。此外，由于当 $\ell$ 充分大时
$\mathcal J^\ell_{\mathrm{reg},K}(T,\{A_\alpha\})$ 在 $\mathcal J^\ell$ 中稠密，仿照
定理~\ref{thm:3-1-6} 的证明可知，
$\mathcal J_{\mathrm{reg},K}(T,\{A_\alpha\})$ 在
$\mathcal J_\tau(M,\omega)$ 中稠密。因此
\[
  \mathcal J_{\mathrm{reg}}(T,\{A_\alpha\})
  =\bigcap_{K\in\mathbb N}
  \mathcal J_{\mathrm{reg},K}(T,\{A_\alpha\})
\]
\sourcefootnote{原书把交集指标印作所有实数 $K>0$，随即称其为“可数交”。这里的
$K$-有界子集随 $K$ 单调并穷竭\BookRecordIndex{T0246}模空间，故取正整数这一可数共尾子集即可，依此改正。}
是开稠密集的可数交，因而是 $\mathcal J_\tau(M,\omega)$ 中的剩余子集。这证明了
定理~\ref{thm:6-2-6} 的 (ii)。
\end{proof}

\begin{exercise}\label{ex:6-2-9}
证明投影
$\pi^\ell:\widetilde{\mathcal M}^*_{0,T}
(\{A_\alpha\};\mathcal J^\ell)\to\mathcal J^\ell$
是一个弗雷德霍姆映射，其指标为
\[
  \mu(A,T)+\dim G_T
  =2n+2c_1(A)+2k+4e(T).
\]
\end{exercise}

与第 3 章一样，为了在\BookRecordIndex{T0184}简单稳定映射的\BookRecordIndex{T0246}模空间之间得到光滑配边，我们引入一类适当的\BookRecordIndex{T0088}殆复结构
\BookRecordIndex{T0407}正则同伦。设 $J_0,J_1\in\mathcal J_{\mathrm{reg}}(T,\{A_\alpha\})$，并设
\[
  [0,1]\longrightarrow\mathcal J_\tau(M,\omega),
  \qquad \lambda\longmapsto J_\lambda
\]
是 $M$ 上被 $\omega$ \BookRecordIndex{T0372}驯服的\BookRecordIndex{T0088}殆复结构的一条光滑同伦。与之相应的\BookRecordIndex{T0184}简单稳定映射\BookRecordIndex{T0246}模空间为
\[
  \mathcal W^*(\{A_\alpha\};\{J_\lambda\})
  :=\bigl\{(\lambda,\mathbf u)\bigm|
  0\le\lambda\le1,\quad
  \mathbf u\in\mathcal M^*(\{A_\alpha\};J_\lambda)\bigr\},
\]
而以 $T$ 为模型的空间为
\[
  \widetilde{\mathcal W}^*_{0,T}(\{A_\alpha\};\{J_\lambda\})
  :=\bigl\{(\lambda,\mathbf u,\mathbf z)\bigm|
  0\le\lambda\le1,
  \ (\mathbf u,\mathbf z)\in
  \widetilde{\mathcal M}^*_{0,T}(\{A_\alpha\};J_\lambda)\bigr\},
\]
其关于\BookRecordIndex{T0413}重参数化群的商记为
\[
  \mathcal W^*_{0,T}(\{A_\alpha\};\{J_\lambda\})
  :=\widetilde{\mathcal W}^*_{0,T}(\{A_\alpha\};\{J_\lambda\})/G_T.
\]
用 $\mathcal J_\tau(J_0,J_1)$ 表示 $M$ 上所有从 $J_0$ 到 $J_1$ 的、被 $\omega$ \BookRecordIndex{T0372}驯服\BookRecordIndex{T0088}殆复
结构的光滑同伦所成的集合。

\begin{definition}\label{def:6-2-10}
若同伦 $\{J_\lambda\}\in\mathcal J_\tau(J_0,J_1)$ 满足
$J_0,J_1\in\mathcal J_{\mathrm{reg}}(T,\{A_\alpha\})$ 以及下列条件，就称它（对
$T$ 和 $\{A_\alpha\}$）\emph{正则}。
\begin{enumerate}[label=(\roman*),leftmargin=3em]
\item 对每个 $\lambda\in[0,1]$、每个元组
$\mathbf u\in\mathcal M^*(\{A_\alpha\};J_\lambda)$ 和每个 $\alpha\in T$，
\[
  \Omega^{0,1}(S^2,u_\alpha^*TM)
  =\operatorname{im}D_{u_\alpha}+\mathbb R v_\lambda,
\]
其中
$v_\lambda:=(\partial_\lambda J_\lambda)du_\alpha\circ j$
是道路 $\lambda\mapsto J_\lambda$ 的切向量的像。
\item \BookRecordIndex{T0283}求值映射
\[
  \mathcal W^*(\{A_\alpha\};\{J_\lambda\})\times Z(T)
  \longrightarrow M^E,
  \qquad
  (\lambda,\mathbf u,\mathbf z)
  \longmapsto\{u_\alpha(z_{\alpha\beta})\}_{\alpha E\beta}
\]
与 $\Delta^E$ 横截。
\end{enumerate}
\end{definition}

与定义~\ref{def:6-2-1} 的情形一样，条件 (i) 蕴含\BookRecordIndex{T0283}求值映射的定义域是流形，故条件 (ii)
有意义。显式地说，这一条件如下：对每个
\[
  (\lambda,\mathbf u,\mathbf z)
  \in\widetilde{\mathcal W}^*_{0,T}
  (\{A_\alpha\};\{J_\lambda\})
\]
以及每族满足
$v_{\alpha\beta}\in T_{u_\alpha(z_{\alpha\beta})}M$、
$v_{\alpha\beta}+v_{\beta\alpha}=0$ 的切向量，都存在向量场
$\xi_\alpha\in\Omega^0(S^2,u_\alpha^*TM)$ 和切向量
$\zeta_{\alpha\beta}\in T_{z_{\alpha\beta}}S^2$，使得
$D_{u_\alpha}\xi_\alpha\in\mathbb Rv_\lambda$，并且对所有 $\alpha E\beta$，
\[
  v_{\alpha\beta}
  =\xi_\alpha(z_{\alpha\beta})
   +du_\alpha(z_{\alpha\beta})\zeta_{\alpha\beta}
   -\xi_\beta(z_{\beta\alpha})
   -du_\beta(z_{\beta\alpha})\zeta_{\beta\alpha}.
\]

记 $\mathcal J_{\mathrm{reg}}(T,\{A_\alpha\};J_0,J_1)$ 为对
$T$ 和 $\{A_\alpha\}$ 正则的同伦集合，并记
$\mathcal J_{\mathrm{reg}}(M,\omega;J_0,J_1)$ 为这些集合对所有树 $T$ 和所有
$\{A_\alpha\}$ 的交。以下结果是定理~\ref{thm:6-2-6} 的直接推广。将
定理~\ref{thm:6-2-6} 的证明与定理~\ref{thm:3-1-8} 的证明相结合即可证明它；细节留给
读者作为习题。

\begin{theorem}\label{thm:6-2-11}
设 $T$ 是一棵 $k$-\BookRecordIndex{T0084}带标号树，固定满足稳定性条件~\eqref{eq:6-1-1} 的映射
\[
  T\longrightarrow H_2(M;\mathbb Z),
  \qquad \alpha\longmapsto A_\alpha,
\]
并设 $J_0,J_1\in\mathcal J_{\mathrm{reg}}(T,\{A_\alpha\})$。
\begin{enumerate}[label=(\roman*),leftmargin=3em]
\item 若
$\{J_\lambda\}\in
\mathcal J_{\mathrm{reg}}(T,\{A_\alpha\};J_0,J_1)$，则\BookRecordIndex{T0246}模空间
$\mathcal W^*_{0,T}(\{A_\alpha\};\{J_\lambda\})$ 是光滑定向带边流形，而且
\[
  \partial\mathcal W^*_{0,T}(\{A_\alpha\};\{J_\lambda\})
  =\mathcal M^*_{0,T}(\{A_\alpha\};J_1)
   \cup\bigl(-\mathcal M^*_{0,T}(\{A_\alpha\};J_0)\bigr).
\]
\item 集合
$\mathcal J_{\mathrm{reg}}(T,\{A_\alpha\};J_0,J_1)$ 是
$\mathcal J_\tau(J_0,J_1)$ 中的\BookRecordIndex{T0300}剩余集。
\end{enumerate}
\end{theorem}
